% HD-F1-T13. GPT-6 Astra Pro, /root/transport; 2026-11-22 UTC.
\section[HD-F1-T13: sporadic source exclusion]{HD-F1-T13: the sporadic projective source gate is empty}
\label{sec:f1-sporadic-source-exclusion}

No source in the remaining moving-rank-at-least-five branch of
T2-OTHER has almost simple projective coefficient image with a
sporadic simple socle. This is an exclusion at the original
coefficient representation and its prescribed primitive Jacobian
corner. It introduces no additional detector family.

Here and below the eight fields are
\(K=\mathbf Q(i)\) and
\(K=\mathbf Q(\sqrt{-d})\),
\(d=2,7,11,19,43,67,163\). Write the actual faithful absolutely
irreducible coefficient representation as
\(G\subset GL_K(V)\), \(\dim_K V=m\). The connected full cover
is \(C\to\mathbf P^1\), with nonidentity inertia matrices
\(g_1,\ldots,g_b\). The original parabolic multiplicity space
has dimension six and period rank nine. Its exact Euler formula
and pointed-base parameter bound give
\[
 c_j=\operatorname{rank}(g_j-1),\qquad
 \sum_{j=1}^{b}c_j=2m+6,\qquad b\ge12.
 \tag{F1T13.1}
\]
In the remaining branch \(c_j\ge5\), so \(m\ge27\).
The proof below actually gives \(m\le15\) for every sporadic
projective image satisfying (F1T13.1), and \(m\le6\) unless
the socle is \(Fi_{24}'\).

\paragraph{1. Conjugate generation controls every eigenspace of the actual lift.}
By absolute irreducibility the centre \(Z=Z(G)\) consists
exactly of scalar matrices. Put \(\overline G=G/Z\), and take
its sporadic socle \(S\), with the actual inclusion
\(S\le\overline G\le\operatorname{Aut}(S)\).
Let \(P\) be the inverse image of \(S\) and set \(L=P'\).
Because \(P/Z=S\) is perfect,
\[
 P=ZP',\qquad
 L=[ZP',ZP']=[P',P'],\qquad P=ZL.
 \tag{F1T13.2}
\]
Thus \(L\) is perfect, normal in \(G\), and a central
extension of \(S\). The invariant space \(V_{\mathbf C}^L\)
is \(G\)-stable. It is zero: otherwise irreducibility would
make it all of \(V_{\mathbf C}\), contradicting faithfulness
and the nontrivial quotient \(L\twoheadrightarrow S\).
No irreducibility of \(V_{\mathbf C}|_L\) has been assumed.

For a noncentral \(g\in G\), let \(x\) be its image in
\(\overline G\). Write \(a=\alpha_S(x)\) for the minimum
number of conjugates of \(x\) generating \(\langle S,x\rangle\).
Lift conjugating elements in \(\langle S,x\rangle\) to
\(G\), obtaining actual conjugates \(g_1',\ldots,g_a'\) of
\(g\), and put \(H=\langle g_1',\ldots,g_a'\rangle\).
Its projective image contains \(S\), so \(H\cap P\) maps
onto \(S\). Consequently
\[
 P=Z(H\cap P),\qquad
 L=P'=(H\cap P)'\subseteq H.
 \tag{F1T13.3}
\]
For any eigenvalue \(\lambda\) of the actual matrix \(g\),
every generator \(g_i'\) acts by the same scalar \(\lambda\)
on
\[
 U_\lambda=\bigcap_{i=1}^{a}\ker(g_i'-\lambda).
\]
This space is \(H\)-stable, and the resulting action on it
is a scalar character of \(H\). The perfect subgroup
\(L\subseteq H\) acts trivially. Thus \(U_\lambda=0\),
by \(V_{\mathbf C}^L=0\). Codimensions of intersections now give
\[
 m\le a\,\operatorname{codim}\ker(g-\lambda),\qquad
 \operatorname{codim}\ker(g-\lambda)\ge m/\alpha_S(x).
 \tag{F1T13.4}
\]
This estimate applies to every eigenvalue, including to central
scalar multiples of lifts of projective involutions. It never
substitutes the order of \(x\) for the order of \(g\).
A nonidentity scalar matrix itself has moving rank \(m\).

\paragraph{2. The published uniform bound and its precise exceptions.}
Revin--Zavarnitsine, \emph{Generation by Conjugate Elements of
Finite Almost Simple Groups With a Sporadic Socle},
\emph{Bull. Irkutsk State Univ. Ser. Math.}~49 (2024),
135--142, Theorem~2, pp.~138--139
(DOI \texttt{10.26516/1997-7670.2024.49.135};
Theorem~3 in arXiv:2402.03673), proves
\[
 \alpha_S(x)\le4
 \quad\text{except}\quad
 (Fi_{22},2A):6,\quad (Fi_{23},2A):6,\quad
 (Fi_{24}',2C):5.
 \tag{F1T13.5}
\]
The theorem includes inner and outer automorphisms of every
order. Its proof and the preceding outer-involution theorem
were read in the published version and compared with the
complete arXiv proof. Thus only the two inner Fischer classes
need an estimate beyond (F1T13.4).

\paragraph{3. Every actual Fischer central lift is covered by the exact tables.}
The ATLAS records the Schur multipliers of \(Fi_{22}\) and
\(Fi_{23}\) as cyclic of order six and trivial, respectively.
The kernel of the perfect central extension
\(L\twoheadrightarrow S\) is a quotient of that multiplier.
It is also a subgroup of the scalar roots of unity in \(K\):
the scalar entries lie in \(K\) because the original matrices
do. For the present fields \(|\mu_K|\) is four or two.
It follows that the actual \(L\) is one of
\[
 Fi_{22},\quad 2.Fi_{22},\quad Fi_{23}.
 \tag{F1T13.6}
\]
In the two exceptional cases the projective element lies in
\(S\), so the actual inertia \(g\in P=ZL\) has the form
\(g=z\ell\), where \(\ell\in L\) lies over class \(2A\)
and \(z\) is a scalar. Multiplication by \(z\) permutes the
eigenvalue labels and preserves their multiplicities.

Here is the finite exact calculation used for \(\ell\).
The ordinary ATLAS character tables in GAP CTblLib have
respectively \(65,114,98\) irreducible rows. Class \(2A\)
is selected by its order and its class size, respectively
\(3510\) and \(31671\). The stored central quotient fusion
for \(2.Fi_{22}\to Fi_{22}\) selects both lifted classes,
numbered \(3,4\), rather than one chosen lift.
For every nontrivial irreducible character \(\chi\) and
every selected element of order \(o\), the multiplicity of
\(\zeta_o^a\) is computed exactly as
\[
 n_a=\frac1o\sum_{k=0}^{o-1}
          \chi(\ell^k)\zeta_o^{-ak}.
 \tag{F1T13.7}
\]
All multiplicities are nonnegative integers and sum to
\(\chi(1)\). The exact minima are
\[
\begin{array}{c|c|c}
 L&\text{lifted classes}&
 \displaystyle\min_{\chi\ne1}
 \frac{\chi(1)-\max_a n_a}{\chi(1)}\\ \hline
 Fi_{22}&2A&11/39\\
 2.Fi_{22}&2B,2C&11/39\\
 Fi_{23}&2A&160/391
\end{array}
 \tag{F1T13.8}
\]
The first minimum occurs in the degree-78 character with
trace \(-34\), giving eigenvalue multiplicities \(22,56\).
Both minima are strictly greater than \(1/4\).

The complete script, exact 387-row output, log and dependency
hashes are saved in
\path{work/root/fischer_review/}.
It checks every pair of character rows for orthogonality,
the degree-square sums and class-size sums. An independent
audit reconstructs all 387 multiplicity pairs as
\((\chi(1)+\chi(\ell),\chi(1)-\chi(\ell))/2\), because
the selected lifted classes all have order two; it verifies
every stored ratio and the three exact minima. The scalar
\(z\) above may still give the original \(g\) order four
or another order. The proof uses the multiplicities, so
that change does not affect (F1T13.8).

The restriction of the actual \(V_{\mathbf C}\) to \(L\)
is a direct sum of nontrivial irreducibles, since
\(V_{\mathbf C}^L=0\). For any fixed eigenvalue, its
multiplicity on each summand is bounded by the largest
multiplicity on that summand. Summing (F1T13.8) therefore
gives codimension at least \(m/4\) for every eigenspace
of the original \(g=z\ell\), with arbitrary constituent
multiplicities. No irreducible summand replaces the
original coefficient representation or its factor.

\paragraph{4. The contradiction at the original source and its framework maps.}
Equations (F1T13.4)--(F1T13.8) give
\(c_j\ge m/5\) for every nonidentity actual inertia.
Hence (F1T13.1) implies
\[
 2m+6\ge\frac{12m}{5},\qquad m\le15.
 \tag{F1T13.9}
\]
Unless \(S=Fi_{24}'\), the stronger \(c_j\ge m/4\)
gives \(m\le6\). Both bounds contradict \(m\ge27\) in
the remaining branch.

The geometric map throughout is the identity of the original
cover \(C\to\mathbf P^1\) and of its middle-extension
coefficient sheaf. Its stalk at a branch value is exactly
\(V^{g_j}\); the alternating Euler dimension of this
same sheaf is (F1T13.1). The period map factors through
the finite Hurwitz cover of the pointed-sphere space of
dimension \(b-3\), which is the source of \(b\ge12\).
Thus the proof counts the original true branch values.
Restriction to \(L\) is used only to bound a matrix's
eigenspaces inside the original \(V\).

In particular the original rational primitive projector
\(e\in\mathbf Q[G]\), the actual image of an integer
multiple of \(e\) in \(JC\), and its restricted theta
polarization are unchanged. The map on the Poincare
biextension, on parabolic cup products and on the selected
six-dimensional Hermitian space is their identity, with
polarization multiplier one. The same identity is used
in the variational Hodge, Grothendieck-group and Adams
realizations already attached to that factor. No separate
factor, character constituent, or virtual class is put in
place of the prescribed corner. The result removes this
source family from the missing gate and makes no claim
to close the other projective-image families.
